\documentclass{article}
\usepackage{icmlko}
\usepackage{graphicx}
\usepackage{booktabs}
\usepackage{amsmath}

\title{「원리상 불가능」이라 적어 둔 문을 열어 보니 열쇠가 옆 파일에 있었다 \\
  {\large 그리고 손으로 센 0 은 러너로도 0 이었다 --- 다만 A 등급 안에서만}}

\begin{document}
\maketitle

\begin{abstract}
직전 노트는 두 문장을 남겼다. 하나는 프레시 티처가 \textbf{손으로 센} 표
(「A 등급 짝 중 날짜형이 아니면서 판 채점 유보를 분모로 쓰는 짝은 0 개다」)이고,
다른 하나는 직전 노트 자신이 적은 \textbf{「팝업은 축 파일이 없어 되짚기가
\emph{원리상 불가능}하다」}였다. 이 노트는 둘을 러너로 다시 잰다.
\claimx{손으로 센 0 은 러너로도 0 이다 --- 그러나 그것은 A 등급에 한정된 사실이다.}
{A 등급 21짝(T2 를 더하면 23짝) 안에서 (형 $\neq$ d) $\wedge$ (분모 = D3) 인 짝이
하나라도 나오면 무너진다. 실측 \textbf{0} 이다. 같은 수가 105짝 전량에서는 \textbf{66} 이다.}
\claimx{그리고 「원리상 불가능」은 거짓이었다 --- 열쇠는 옆 파일에 있었다.}
{판 채점 유보 팝업 키 중 원천 레코드에 \emph{못 붙는} 키가 하나라도 있으면 무너진다.
실측 \textbf{0} 이다: 유보 \textbf{65} = 원천 \textbf{59} + 시장 \textbf{6} + 못 붙음 \textbf{0}.}
$\Rightarrow$ 이 노트가 뒤집는 것은 자료가 아니라 \textbf{저장소가 자기에 대해 적어 둔 문장}이다.
\end{abstract}

\section{두 문장, 하나의 병}

직전 사이클의 최대 발견은 그 사이클이 쓴 것이 아니었다.
프레시 티처가 직전 산출물의 판정 표를 \textbf{손으로} 전수 세어
「A 21 = 팝업 12(전부 형 c $\cdot$ 전부 분모 D1) + 날짜 9(전부 형 d $\cdot$ 전부 분모 D3)」
이고 그 교집합이 \textbf{0} 이라고 적었다.
이 저장소의 규약대로면 \textbf{손으로 센 수는 후보이지 판정이 아니다.}

같은 사이클이 별도로 적은 문장이 하나 더 있다 ---
\textbf{「팝업은 축이 \texttt{npz} 라 D3 되짚기가 원리상 불가능하다」.}
둘은 서로 다른 자리에 있지만 병이 같다.
\textbf{하나는 세지 않고 적었고, 하나는 열어 보지 않고 적었다.}

\section{재기 전에 자백한다}

사전등록 \texttt{docs/prereg\_904\_xtab.md} 는 커밋 \texttt{21c272113}
(2026-08-11T07:27:57+09:00)로 들어갔고, 측정 산출물 \texttt{runners/out904\_xtab.json} 의
스탬프는 \textbf{2026-08-11T07:32:43 $\to$ 2026-08-11T07:32:49}(5.434 초)다.
\texttt{stamp} 은 \textbf{실행 시작}에서 부른다. 커밋 시각이 측정 시작보다 앞선다.

🔴 그런데 이 사전등록에는 \S 6(판정 어법)보다 \textbf{먼저 오는 \S 0} 이 있다.
직전 노트가 *「수가 강제했다」*라고 적었는데 그 결론은 측정 전에 이미 문서에 있었다 ---
그래서 이 노트는 반대로 한다. \textbf{사전등록을 쓰기 위해 무엇을 먼저 열어 봤는지를
표로 자백하고 시작한다}: \texttt{popup\_v2.npz} 의 키 여섯,
\texttt{popup\_v2\_meta.json} 의 길이와 필드, 원천 레코드의 \texttt{record\_id},
그리고 챔피언 팝업 행 89 $\cdot$ 유보 65.

\claimx{그러므로 「팝업 축에 열쇠가 있다」는 이 사이클의 \emph{발견}이 아니라 \S 0 의 \emph{자백}이다.}
{이 논문$\cdot$산출물$\cdot$커밋 제목 어디에도 「우리가 찾아냈다」가 나오면 무너진다.}
러너가 한 일은 그 사실에 \textbf{수와 분모를 단 것}이다.
그리고 \S 0 이 \emph{아직 안 본 것}으로 남겨 둔 다섯 가지가 이 노트의 진짜 예측이다.

\section{배선 검사 --- 이번엔 게이트의 복제본을 안 때린다}

직전 노트의 배선 검사 넷 중 \textbf{셋이 게이트가 아니라 게이트의 복제본}을 시험했다.
생산 코드의 \texttt{assert} 를 지워도 「4/4 발화 $\cdot$ 통과 True」가 그대로 나왔다.
그래서 이번 검사는 셋을 지킨다: \textbf{(1)} 검사가 생산 함수를 \emph{부른다} ---
증거로 함수 안에 \textbf{진입 계수기}를 두고 검사 전후 증가분을 산출물에 싣는다,
\textbf{(2)} 음성 대조(정상 입력에서 발화 안 함),
\textbf{(3)} \textbf{「심은 결함을 지우면 검사가 통과로 바뀌는가」}를 검사마다 실제로 돌린다.

\begin{center}
\begin{tabular}{llcc}
\toprule
검사 & 부른 생산 함수 & 심으면 발화 & 지우면 통과 \\
\midrule
형 결측 & \texttt{pair\_rows} & ✓ & ✓ \\
분모 딱지 넷 밖 & \texttt{crosstab} & ✓ & ✓ \\
축 정렬 밀기 & \texttt{popup\_keys} & ✓ & ✓ \\
키 접두 벗기기 & \texttt{resolve\_popup\_d3} & ✓ & ✓ \\
\bottomrule
\end{tabular}
\end{center}

심은 결함 \textbf{4} $\cdot$ 발화 \textbf{4} $\cdot$ 음성 대조 통과 \textbf{4} $\cdot$
지우면 통과로 바뀐 수 \textbf{4} $\cdot$ 진입 계수기가 실제로 늘어난 검사 \textbf{4}.
🔴 마지막 열이 직전 노트에 없던 열이다.

\section{(가) 교차표 --- 0 은 재현된다. 그리고 A 등급 밖에선 0 이 아니다}

\begin{figure}[h]\centering
\includegraphics[width=\linewidth]{figs/fig1.pdf}
\caption{(a) A$\wedge$T1 21짝의 형 $\times$ 분모. 괄호는 105짝 전량.
형 c 는 전부 D1(12), 형 d 는 전부 D3(9) --- \textbf{두 칸만 차 있다}.
(b) 같은 수를 네 모집단에서. 붉은 막대가 0 이다.}
\end{figure}

\begin{center}
\begin{tabular}{lrrr}
\toprule
모집단 & 짝 수 & (형 $\neq$ d) $\wedge$ D3 & 형 $\times$ 분모 \\
\midrule
A $\wedge$ T1 & 21 & \textbf{0} & c$\times$D1 12 $\cdot$ d$\times$D3 9 \\
A 전량(T2 포함) & 23 & \textbf{0} & c$\times$D1 14 $\cdot$ d$\times$D3 9 \\
T1 전량 & 81 & 47 & --- \\
105짝 전량 & 105 & \textbf{66} & --- \\
\bottomrule
\end{tabular}
\end{center}

\claimx{손으로 센 0 이 러너로 재현됐다 --- 그리고 분모를 바꿔도 0 이다.}
{A$\wedge$T1 21 에서든 A 전량 23 에서든 1 이상이 나오면 무너진다.}
형과 분모가 겹치는 것은 처방서의 오타가 아니라 \textbf{자료의 구조}다.
「날짜형을 빼고 D3 위에서 효과를 재라」는 처방은 \textbf{고를 수 있는 짝이 하나도 없다.}

🔴 그러나 같은 수가 105짝 전량에서는 \textbf{66}(등급 B 51 $\cdot$ C 15)이다.
\textbf{0 은 A 등급에 한정된 사실이지 자료 전체의 사실이 아니다.}
이 문장은 직전 사이클의 티처가 안 적었고, 적었으면 처방이 달라졌을 문장이다.

\section{(나) 축 파일 --- 열쇠는 옆 파일에 있었다}

\begin{figure}[h]\centering
\includegraphics[width=\linewidth]{figs/fig2.pdf}
\caption{(a) 판 채점 유보 팝업 65 키의 되짚기 --- 못 붙은 키 0.
(b) 12 A$\wedge$T1 팝업 짝의 비결측(분모 D1 = 380) 대 D3 덮음.}
\end{figure}

\texttt{data/state/popup\_v2.npz} 는 키 여섯을 갖는데 식별자로 쓸 수 있는 열이
\textbf{0} 개다 --- 여기까지는 직전 노트가 맞다. 그런데 같은 폴더의
\texttt{popup\_v2\_meta.json} 이 \textbf{길이 376 의 리스트}이고
\texttt{npz} 의 행 수와 같으며 각 원소가 \texttt{id} 를 갖는다.
그 \texttt{id} 는 \texttt{data/records/*.json} 의 \texttt{record\_id} 와 \emph{같은 이름}이다.

키를 복원할 때 \textbf{생산 코드와 원소별로 댔다}. 챔피언 배선은
\texttt{ff753.shell(base())} $\to$ \texttt{lab.loop.\_idol(wide\_pop="grades")} $\to$
\texttt{lab.popupset.build} 이고, 이 경로가 실제로 낸 \texttt{t}(분수 날짜) $\cdot$
\texttt{y}(일평균 방문) $\cdot$ 행 수 $\cdot$ \texttt{counting} 도수 넷과
어긋난 원소가 \textbf{0} 이다. 좁은 판(\texttt{state.slots.load\_popup} $\cdot$ 75행)의
\texttt{groups} $\cdot$ \texttt{times} 와도 어긋남 \textbf{0} 이다 --- \textbf{두 생산 경로를 따로 댔다.}

\begin{center}
\begin{tabular}{lr}
\toprule
분모 & 값 \\
\midrule
\texttt{popup\_v2.npz} 행 = \texttt{popup\_v2\_meta.json} 길이 & 376 \\
챔피언 팝업 도메인 행 (D4) & 89 \\
챔피언 팝업 판 채점 유보 (D3) & \textbf{65} \\
\quad $\hookrightarrow$ \texttt{data/records} 에 붙은 수 & 59 \\
\quad $\hookrightarrow$ \texttt{data/market\_records} 에 붙은 수 & 6 \\
\quad $\hookrightarrow$ 🔴 못 붙은 수 & \textbf{0} \\
\bottomrule
\end{tabular}
\end{center}

\claimx{「원리상 불가능」이 아니라 「안 써 봤다」였다.}
{못 붙은 키가 1 이상이면 무너진다. 실측 0.}

그래서 직전 노트가 12번 반복해 적은 \textbf{「못 셌다 --- 축 파일이 없다」}가
처음으로 수가 됐다. A$\wedge$T1 팝업 12짝의 \textbf{D3 덮음}은
최소 \textbf{13}(\texttt{area\_sqm}) 에서 최대 \textbf{58}(\texttt{entry\_friction} 등 다섯)이고,
\textbf{덮음이 비결측과 같은 짝은 0 개}다.
직전 노트가 D1 위에서 짝 \textbf{95} 로 잰 그 열은 판 유보 위에서는 \textbf{58} 이다.

⚠ \textbf{두 수를 잇지 않는다}(분모가 다르다). 그리고 이 노트는
\textbf{효과를 하나도 재지 않았다} --- $\delta$ 도 구간도 없다.
여기까지가 「분모를 옮길 수 있는가」이고, 「옮기면 무엇이 나오는가」는 안 쟀다.

\section{맺음 --- 다음이 무엇인지는 이 노트가 안 적는다}

직전 노트는 판정문에 「다음은 표본으로 간다 --- 이 수가 강제한다」를 적었고,
그 문장은 측정 전에 이미 사전등록에 있었다. 그래서 이 노트의 사전등록 \S 6 은
\textbf{세 번째 칸}을 미리 못 박았다:

\claimx{이 사이클은 「다음은 X 로 간다」를 판정문에 안 쓴다.}
{판정문 어디에 다음 사이클의 방향이 적히면 무너진다.}

낼 수 있는 것은 두 수뿐이다 --- 교차표의 \textbf{0}(그리고 그 밖의 \textbf{66}),
그리고 되짚기의 \textbf{65 = 59 + 6 + 0}.
방향은 다음 티처와 다음 사이클의 \S ⓪ 이 정한다.

\textbf{못 잰 것}(「없다」가 아니다): 효과 크기 $\cdot$ 아이돌의 되짚기(안 봤다) $\cdot$
\texttt{brand\_key}$\cdot$\texttt{space\_key} 로의 이름 붙이기(안 써 봤다) $\cdot$
날짜형 아홉 짝 $\cdot$ 시장팝업으로 붙은 6 키의 분모 재계산.

\end{document}
